\documentclass[10pt]{beamer}

\usetheme{Madrid}
\usefonttheme{professionalfonts}
\setbeamertemplate{footline}[frame number]
\setbeamercovered{transparent}

\usepackage{algorithm}
\usepackage{algpseudocode}
\usepackage{xcolor}
\renewcommand{\algorithmiccomment}[1]{\hfill$\triangleright$\textcolor{WildStrawberry}{#1}}
\usepackage{float}

\begin{document}

\title{Gurobi}
\author{Polina Evgrafova,Catalin Viscun}
\date{November 18, 2025}

\begin{frame}
    \maketitle
\end{frame}
\begin{frame}{About a license}
    \begin{enumerate}\Large\textbf{How to install gurobi}\normalsize
        \item Create an academic account on a official site, you'll need to type an institution by yourself.
        \item On the "My account" page choose subpage "Licenses" and request "Named-User Academic"
        \item For python use:
        \\ \texttt{pip install gurobipy}
        \\ \texttt{grbgetkey key-you'll-get-after-obtaining-a-license}
        \item For the purpose of wide usage of a optimization application we should probably use a "Educational Institution Site" license
    \end{enumerate}
    
\end{frame}



\begin{frame}{Heuristics Overview}
    \begin{itemize}\Large\textbf{What heuristics achieve}\normalsize
        \item Provide good feasible solutions \textbf{without guaranteeing optimality}
        \item \textbf{Reduce} computation time for large or hard MIP instances
        \item \textbf{Improve} lower bounds via feasible incumbents
        \item Allow \textbf{hybrid} workflows: exact MILP + heuristic search
    \end{itemize}

    \bigskip
    \begin{itemize}\Large\textbf{Gurobi's role}\normalsize
        \item Gurobi implements a suite of primal heuristics automatically
        \item Heuristics can be strengthened through solver parameters
        \item \textbf{Custom heuristics} can be executed by modifying variables, fixing subsets, or adding constraints
    \end{itemize}
\end{frame}

\begin{frame}{Built-in Heuristics}
    \begin{itemize}\Large\textbf{Frequently applied primal heuristics}\normalsize
        \item \textbf{Feasibility Pump:} rapidly constructs first feasible solutions
        \item \textbf{RINS:} neighborhood search based on LP relaxation and incumbent
        \item \textbf{Diving Heuristics:} progressively fix variables based on LP values
        \item \textbf{Local Search:} improvement steps on incumbent solutions
        \item \textbf{Rounding:} simple rounding from LP relaxation
        \item \textbf{MIP Polishing:} solution improvement in the final phase
    \end{itemize}

    \bigskip
    \begin{itemize}\Large\textbf{Benefits}\normalsize
        \item Strong feasible solutions early in the search
        \item Better MIP gap reduction
        \item Reduced branch-and-bound tree growth
    \end{itemize}
\end{frame}

\begin{frame}{Custom Heuristics}
    \Large\textbf{Strategies that integrate well with Gurobi}\normalsize
    \begin{itemize}
        \item \textbf{LP Rounding:} solve LP relaxation and round fractional values
        \item \textbf{Relax-and-Fix:} iteratively fix subsets of variables
        \item \textbf{Fix-and-Dive:} fix easy variables first and dive into solutions
        \item \textbf{Local Branching:} explore solutions within a defined Hamming distance
        \item \textbf{Large Neighborhood Search:} fix a neighborhood and re-optimise
    \end{itemize}

    \bigskip
    \begin{enumerate}\Large\textbf{Workflow}\normalsize
        \item Construct relaxed or reduced subproblem
        \item Solve subproblem with Gurobi (fast exact solver)
        \item Extract feasible solution and use as warm start
        \item Repeat to explore new neighborhoods
    \end{enumerate}

    \bigskip
    \textbf{Reference implementation:} \url{https://www.gurobi.com/documentation/}
\end{frame}

\begin{frame}{Linear integer programming}
    \begin{itemize}\Large\textbf{What are we able to do}\normalsize
        \item Introduce variables for constrains and auxiliary variables for computing
        \item Add \textbf{hard constrains} using \texttt{addConstrs} for simple constrains and further \texttt{addConstrs-} for composite ones
        \item Add \textbf{soft constrains} using by introducing an additional variable equal to 1 when we allow to break a constraint
        \item Introduce a \textbf{minimisation} function by \texttt{setObjective}
        \item Compute parallel
        \item Set time limitation
    \end{itemize}
    \bigskip 
    \textbf{Where to find more:} \url{https://docs.gurobi.com/projects/optimizer/en/current/reference/python/model.html}
\end{frame}

\begin{frame}{Linear integer programming}
\begin{itemize}\Large\textbf{What we need to implement}\normalsize
    \item Model constrains
    \item A way to collect and store constrains
    \item A way to correlate weights with soft constrains
    \item Tests to determine well suited weights
\end{itemize}
    
\end{frame}


        
\end{document}